As enchentes estão entre os desastres naturais mais frequentes e impactantes do mundo, afetando milhões de pessoas e
causando prejuízos econômicos e sociais significativos todos os anos. Esses eventos, caracterizados pelo transbordamento
de corpos d’água ou pelo acúmulo excessivo de água em áreas urbanas e rurais, são resultado de uma complexa interação
entre fatores naturais e atividades humanas. No contexto das mudanças climáticas, a frequência e a intensidade das
enchentes têm aumentado em diversas regiões, tornando-se um desafio global que exige atenção urgente.

Neste capítulo, será abordado as principais características das enchentes, com foco especial nas enchentes fluviais, que
ocorrem quando rios transbordam. Como estudo de caso destes eventos, serão trazidos eventos recentes de enchentes
catastróficas e os seus impactos, um exemplo emblemático de como a combinação de mudanças climáticas, alterações na
morfologia dos rios e vulnerabilidade urbana pode levar a desastres de grandes proporções. Além disso, serão explorados
os processos hidrodinâmicos que influenciam a ocorrência e a duração das enchentes, bem como a intrínseca
vulnerabilidade das áreas urbanas a esses eventos extremos.

% ######################################################################################################################
% IPCC AR6 WGI - Chapter 8
\subsection{Enchentes e suas características}

\begin{figure}[htb] 
	\caption{\label{cicle} Ciclo da água}
	\begin{center}
	    \includegraphics[scale=0.165]{img/WaterCicle (copy).png}
	\end{center}
	\legend{Fonte: \textit{United States Geological Survey} \cite{usgs}, traduzido pelo autor}
\end{figure}

Enchentes são caracterizadas pelo aumento do nível da água em áreas que normalmente permanecem secas. Esse acúmulo
prolongado ocorre devido à dificuldade de escoamento, seja pela incapacidade do solo de absorvê-la ou pela lentidão na
drenagem natural, como o fluxo em grandes corpos d'água. Em geral, elas podem ser classificadas em três tipos
principais: enchentes pluviais, resultantes de chuvas intensas que excedem a capacidade de absorção do solo, causando
alagamentos; enchentes fluviais, provocadas pelo transbordamento de corpos d'água – geralmente rios – quando o nível da
água ultrapassa suas margens, frequentemente devido a chuvas persistentes, derretimento de gelo ou acúmulo de
sedimentos; e enchentes costeiras, que ocorrem em áreas litorâneas devido ao aumento do nível do mar, grandes ondas ou
fenômenos como ciclones e tsunamis.

As causas de uma enchente nem sempre podem ser facilmente atribuídas a um único evento, pois geralmente dependem de uma
combinação de fatores. No caso das enchentes fluviais, por exemplo, não é apenas a ocorrência de chuvas intensas no
corpo d'água principal ou em seus tributários que provoca o transbordamento. Se, por exemplo, o corpo hídrico tiver
capacidade de escoar a água de forma eficiente, o alagamento pode ser evitado. Muitos corpos d'água, cuja principal
fonte de abastecimento são enchentes pluviais, possuem uma morfologia adaptada para lidar com essa condição. Outros
fatores que contribuem para o transbordamento incluem a área de escoamento disponível em canais naturais e artificiais,
a umidade prévia do solo, sua capacidade de percolação, a presença de ventos fortes e sua interação com o fluxo da água,
a taxa de evaporação após o evento, a capacidade de absorção da vegetação local, o declive do terreno em relação à área
de escoamento, entre outros. Assim, não se pode associar exclusivamente o aumento das chuvas intensas (\autoref{rain})
ao crescimento destes eventos extremos \cite{ipcc_ar6_wg2_chap11}.

Mesmo assim, as enchentes são um fenômeno natural e frequentemente característico de certas regiões. Áreas com invernos
rigorosos, marcados pela acumulação de neve, costumam enfrentar enchentes durante o verão \cite{atashi}, quando o rápido
derretimento do gelo transporta grandes volumes de água de uma região para outra, sustentando ecossistemas que dependem
desse influxo sazonal. Regiões influenciadas por monções — ventos sazonais que trazem para o interior um acúmulo de
umidade e precipitação marítima — também são marcadas por cheias em seus corpos d'água, essenciais para a agricultura
local, como ocorre nos tributários dos rios Ganges e Amarelo. De maneira semelhante, as cheias anuais regulares do Rio
Nilo são fundamentais para a fertilidade da região. 

Entretanto, mesmo sendo uma parte integral do ciclo da água, a presença de assentamentos humanos geralmente corresponde
com a existência de grandes corpos d'água que podem ser, dadas circunstâncias, vulneráveis à enchentes. Essas imundações
podem invadir esses locais de forma inesperada e veloz, colocando vidas em risco, comprometendo a segurança das
comunidades e causando danos estruturais que amplificam a disrupção das atividades humanas. Fatalidades causadas por
afogamentos são comuns em regiões alagadas, não apenas pela rapidez com que as enchentes ocorrem, mas também devido às
características do escoamento das águas, que podem gerar correntes fortes e imprevisíveis. As enchentes também criam
condições propícias para o agravamento de outras condições sanitárias. Elas podem transmitir doenças por meio da água
contaminada, além de comprometer sistemas críticos, como a contaminação de fontes de água potável, a interrupção do
fornecimento de energia elétrica, afetando hospitais e serviços essenciais, e a destruição de cadeias de suprimentos.
Esses fatores podem isolar comunidades, mesmo aquelas não diretamente afetadas pela inundação, dificultando o acesso a
mantimentos e recursos básicos.

As enchentes urbanas representam um dos principais vetores de risco em áreas densamente povoadas, combinando
vulnerabilidade socioambiental com impactos econômicos e humanos catastróficos. 

% Em Roma, por exemplo, exitem registros de cheias no Rio Tibre que remontam ao século 6 BCE, quando a construção da
% Cloaca Máxima marcou um dos primeiros sistemas organizados de gestão hidráulica urbana. Já em 1870, uma falha
% estrutural no canal do rio resultou em uma inundação de 17m acima do nível normal, submergindo monumentos icônicos
% como o Panteão, localizado nos Campos de Marte, uma região considerada como várzea.

% % Photo: Miss I. B. Trevenen, Rome. The Floods in Rome, the Pantheon during the Inundation. Illustration for The %
% Illustrated London News, 15 December 1900. English Photographer (20th Century)

% \begin{figure}[htb] \caption{\label{mars} Enchente de 1900 na entrada do Panteão} \begin{center}
%   \includegraphics[scale=0.55]{img/pantheon.jpg} \end{center} \legend{Fonte: \textit{The Illustrated London News}
%   \cite{london}} \end{figure}

Mais recentemente, destaca-se a enchente de 2003 na provincia de Santa Fe, na Argentina. Depois de dias de chuvas
intensas, o rio Salado transbordou no dia 28 de Abril. Aproximadamente 100 mil pessoas tiveram de ser evacuadas da
região e 160 faleceram. As perdas estimadas da região ficaram na faixa de 500 milhões de pesos argentinos
\cite{santa_fe}.

% Os efeitos das enchentes, no entanto, não se limitam às comunidades urbanas. Atividades econômicas fundamentais, como
% a agricultura, estão entre as mais prejudicadas por esses desastres. A necessidade de grandes extensões de terra e a
% proximidade de corpos d’água para irrigação tornam o setor agrícola particularmente vulnerável. Em situações de
% enchente, a safra pode ser destruída não somente pela força da água, mas também pela permanência prolongada de água
% nessas regiões, já que muitas espécies de plantas não conseguem sobreviver a tais condições. Além da destruição de
% colheitas e da perda de gado, as enchentes frequentemente resultam em danos à infraestrutura essencial, como sistemas
% de irrigação, equipamentos, instalações de processamento e transporte, além de alterações nas terras cultiváveis. De
% acordo com a Organização de Alimentos e Agricultura das Nações Unidas \cite{onu}, enchentes são responsáveis por mais
% da metade de todas as perdas em colheitas causadas por desastres naturais. Entre 2008 e 2018, essas perdas em países
% em desenvolvimento ultrapassaram 21 bilhões de dólares, tornando as enchentes a segunda maior causa de danos
% econômicos no setor agrícola, atrás apenas das secas. Em países em desenvolvimento, onde setores primários têm um peso
% significativo na economia, os efeitos desses desastres são amplamente sentidos de forma indireta no enfraquecimento
% econômico, podendo levar a uma maior situação de vulnerabilidade e despreparo.

% O setor agrícola de Bangladesh, onde 60\% da população depende diretamente da agricultura, é recurrentemente afetado
% por enchentes. Em 2020, inundações no nordeste do país destruíram 220.000 hectares de plantações de arroz, afetando
% 4,7 milhões de pessoas. A FAO estima que, entre 2015 e 2022, perdas médias anuais chegaram a US\$ 1,2 bilhão,
% equivalentes a 2,3\% do PIB agrícola nacional. Esses eventos não apenas comprometem a segurança alimentar, mas também
% perpetuam ciclos de pobreza, já que pequenos produtores perdem até 80\% de sua renda sazonal.

% ######################################################################################################################
\subsection{Tendências de enchentes devido à mudanças climáticas}
De maneira direta e indireta, as atividades humanas alteram profundamente o ciclo hidrológico. Conforme demonstrado na
\autoref{rain}, o aquecimento global não apenas eleva as temperaturas, mas também intensifica a frequência e a magnitude
das chuvas, impulsionado pela maior quantidade de umidade presente na atmosfera. Embora as variações naturais possam
dificultar a identificação de mudanças locais nos eventos de chuva, há evidências de que, quando padrões meteorológicos
naturais promovem chuvas prolongadas em um clima aquecido, a intensidade desses episódios é amplificada pela abundância
de vapor d’água. Isso resulta em um aumento de eventos extremos, mesmo que regionalmente as mudanças possam se
manifestar de formas diversas – com áreas áridas tornando-se ainda mais secas e regiões úmidas sendo afetadas por chuvas
mais intensas.

Entretanto, é importante destacar que chuvas mais intensas nem sempre se traduzem automaticamente em enchentes mais
graves, pois os impactos dependem de múltiplos fatores, como as características da bacia hidrográfica, a topografia, a
duração do evento e o estado de umidade do solo antes da ocorrência da chuva. Além disso, a gestão do uso do solo – como
o desmatamento para expansão agrícola ou a urbanização desordenada – pode acelerar o fluxo de água, direcionando-a de
forma mais rápida para áreas vulneráveis. Em contrapartida, a extração intensiva de água dos rios pode, em certos
contextos, reduzir os níveis dos cursos d’água e, temporariamente, mitigar o risco de enchentes.

Portanto, mesmo que as variações naturais dificultem atribuir as mudanças locais unicamente ao aquecimento global,
existe uma tendência de que o aumento da concentração de gases de efeito estufa contribuem para chuvas diárias cada vez
mais intensas e eventos extremos. Esse cenário, aliado às alterações no manejo do solo e das bacias hidrográficas, leva
à enchentes mais severas em muitas regiões.

Na região Sudeste do Brasil, as tendências de enchentes fluviais se mostram de acordo com os padrões globais. Como é uma
área caracterizada por longos períodos úmidos, espera-se que esses períodos se tornem ainda mais frequentes e intensos.
Em regiões com alta modificação antropogênica, especialmente devido à atividade agrícola, esse agravamento tende a ser
ainda mais pronunciado \cite{chagas}. Eventos de enchentes catastróficas, tradicionalmente calculados para ocorrer uma
vez por século, já apresentam indícios de um período de retorno menor, levando à eventos recordes mais frequentes. De
acordo com a análise de extremos hidrológicos de Vieira et al. (2022) \cite{vieira}, a probabilidade de ocorrência de
enchentes devido à precipitação e fluxo extremo na Bacia Hidrográfica do Guaíba é maior e pode ultrapaçar marcos
históricos, como ocorreu com as enchentes de 2024.

% ######################################################################################################################
\subsection{Morfologia de corpos fluviais e Hidrodinâmica de enchentes}
\label{morf}
Um rio é uma formação hidrológica que transporta água e sedimentos de regiões mais elevadas, conhecidas como cabeceiras,
para áreas de menor altitude em direção à corpos maiores de água, seguindo um caminho definido, denominado canal. Esses
canais desempenham um papel fundamental no ciclo hidrológico, conduzindo a maior parte da água das regiões interioranas
até as áreas costeiras. Ao longo de seu trajeto, os rios erodem o solo e as rochas, transportando sedimentos que,
eventualmente, são depositados em lagos e oceanos \cite{holden}.

% \begin{figure}[htb] \caption{\label{guaiba} Rio Guaíba} \begin{center}
%   \includegraphics[scale=0.113]{img/IMG_20250406_163317.jpg} \end{center} \legend{Fonte: Reprodução própria}
%   \end{figure}

A formação dos canais fluviais depende principalmente de fatores como precipitação, derretimento de neve e água
subterrânea. A precipitação, em particular, destaca-se como um fator crucial. Além de contribuir diretamente para o
abastecimento dos rios, ela promove a infiltração de água no solo, originando aquíferos e água subterrânea. Grande parte
dessa água subterrânea acaba integrando a bacia hidrográfica, também conhecido como sistema de drenagem, por meio de
riachos e córregos, que se combinam para formar canais maiores. Esse processo inicia-se com o escoamento superficial,
onde a água da chuva flui inicialmente como uma fina camada sobre o solo, que se concentra em áreas de maior declive e
começa um processo de erosão do solo, formando pequenos sulcos. Com o tempo, esses sulcos se aprofundam e se fundem,
dando origem a canais mais amplos que estabelecem uma rede de drenagem, formando afluentes que se unem a um rio
principal, criando-se uma complexa rede de drenagem \cite{nelson}.

As bacias hidrográficas podem assumir diversas formas, dependendo das características geográficas e geológicas da região
onde a água escoa. O principal fator é o tipo de rocha presente na área, pois diferentes rochas apresentam resistências
variadas à erosão. Além disso, a topografia local, como a inclinação do terreno e a presença de falhas geológicas,
também ajudam a definir esses sistemas. Dessa forma, o padrão de drenagem de uma região pode ser um indicativo de sua
geologia e topografia subjacentes. Na \autoref{drain}, são apresentados alguns dos tipos mais comuns de bacias
hidrográficas, cada um associado a características específicas do terreno.

\begin{figure}[htb]
	\caption{\label{drain} Exemplos de formas de bacias hidrográficas}
	\begin{center}
	    \includegraphics[scale=0.60]{img/RiverShapes.png}
	\end{center}
	\legend{Fonte: Reprodução própria}
\end{figure}

Na \autoref{drain}.1, observa-se o sistema de drenagem dendrítico, cujo padrão se assemelha aos ramos de uma árvore. É
típico de áreas compostas por granitos ou arenitos, onde a resistência à erosão é uniforme, permitindo o escoamento em
todas as direções, geralmente em planícies ou relevos suaves. A \autoref{drain}.2 mostra o sistema radial, em que os
cursos d’água divergem de um ponto elevado central, comum em regiões montanhosas com forte variação altimétrica. Já a
\autoref{drain}.3 exibe o padrão retangular, com canais angulados, comuns em áreas com falhas ou fraturas na rocha, onde
os rios seguem zonas com diferentes resistências à erosão.

% Esses padrões de drenagem influenciam diretamente a morfologia e a dinâmica dos rios ao longo de seu curso. Nas áreas
% iniciais do canal, próximas à nascente, o gradiente é mais acentuado, refletindo a maior elevação e a forte tendência
% da água de escoar em uma direção específica. À medida que o rio avança em direção à foz, o gradiente diminui, e o
% fluxo se torna mais lento. Nessas regiões de baixo gradiente, é comum a deposição de sedimentos que leva à formação de
% várzeas, ecossistemas úmidos alimentados pelas inundações periódicas do rio, e meandros, curvas sinuosas que resultam
% da erosão nas margens côncavas e da deposição de sedimentos nas margens convexas. Meandros mais acentuados são
% característicos de rios mais velhos e bem estabelecidos. As várzeas, por sua vez, podem servir como uma barreira para
% as enchentes, dado que a sua disposição plana situada adjacentemente aos rios favorece a absorção do excesso de água
% durante cheias, mas também são vulneráveis a inundações prolongadas quando a capacidade de infiltração do solo é
% ultrapassada.

% Essa evolução longitudinal do rio também se manifesta na sua geometria, especialmente na seção transversal do canal. A
% seção transversal de um rio é definida como a área por onde ocorre o fluxo de água no canal. Essa área varia ao longo
% do curso do rio e reflete as características do fluxo, como velocidade, descarga e capacidade de transporte de
% sedimentos. As partes mais profundas do canal geralmente estão associadas às regiões onde a velocidade da corrente é
% maior, pois a água concentra sua energia em áreas específicas, promovendo maior erosão. Além disso, tanto a largura
% quanto a profundidade da seção transversal aumentam conforme a descarga do rio, já que um maior volume de água exige
% uma área mais ampla para escoar. Por fim, a forma da seção transversal costuma ser irregular e variável ao longo do
% trajeto, influenciada por fatores como a composição das rochas, a dinâmica de erosão e deposição de sedimentos, e a
% presença de obstáculos naturais ou artificiais. A segunda parte da \autoref{sinus} ilustra a seção transversal de um
% rio, destacando em especial a relação entre a forma do canal e a velocidade da água.

Os padrões de drenagem condicionam a evolução longitudinal dos rios. Nas cabeceiras, o elevado gradiente promove fluxos
rápidos e erosão. À medida que o gradiente diminui em direção à foz, a redução da velocidade favorece a deposição de
sedimentos, formando várzeas - planícies de inundação adjacentes que atuam como áreas de amortecimento durante cheias -
e meandros por erosão seletiva em margens côncavas e deposição nas convexas. Esta morfologia reflete-se igualmente na
seção transversal do canal, onde a geometria (largura, profundidade e rugosidade) determina a relação entre velocidade,
descarga e capacidade de transporte sedimentar. Como ilustrado na \autoref{sinus}, áreas de maior profundidade
correlacionam-se com velocidades elevadas e erosão, enquanto aumentos na descarga ampliam as dimensões do canal. A forma
transversal é dinamicamente modulada por fatores litológicos, processos erosivo-deposicionais e intervenções antrópicas.

\begin{figure}[htb]
	\caption{\label{floodplain} Exemplo de regiões de várzea inundadas no Rio Severn, Inglaterra}
	\begin{center}
	    \includegraphics[scale=0.35]{img/floodplain.jpg}
	\end{center}
	\legend{Fonte: Tom Blockley \cite{varzea}}
\end{figure}

A seção transversal do canal regula a relação entre largura, profundidade e velocidade do fluxo. Canais estreitos e
profundos concentram a água em alta velocidade, aumentando a capacidade de transporte de sedimentos, mas também elevando
o risco de transbordamento rápido caso a descarga exceda a capacidade de contenção. Em contraste, canais largos e rasos
dissipam parte da energia do fluxo, reduzindo a velocidade, porém exigindo maior área para acomodar volumes elevados de
água. Alterações na seção transversal, como estreitamentos naturais ou acumulo de sedimentos, podem amplificar a
profundidade do fluxo e a pressão sobre as margens, favorecendo transbordamentos.

Além disso, a velocidade da água influencia não apenas a erosão e a deposição, mas também a capacidade do rio de
transportar materiais em suspensão. Irregularidades no leito e nas margens aumentam a resistência ao fluxo, reduzindo a
velocidade da água próxima às margens e intensificando a erosão nessas áreas. Em contraste, a água no centro do canal é
menos afetada por esse efeito, mantendo uma velocidade mais elevada. Além disso, os efeitos da fricção são mais
pronunciados em canais largos e rasos, enquanto em canais estreitos e profundos a resistência é menor. Outro fator
crucial é o gradiente do canal: regiões com maior declive tendem a apresentar águas mais velozes, enquanto áreas de
baixo gradiente favorecem fluxos mais lentos.

A sinuosidade do canal também desempenha um papel importante na distribuição da velocidade da água. Em canais retos, a
velocidade máxima ocorre no centro do canal, devido à menor influência da fricção com as margens. Já em canais sinuosos,
a velocidade é mais alta na parte externa das curvas, onde a força centrífuga empurra a água contra as margens,
resultando em maior profundidade e erosão nesses pontos. Por outro lado, na parte interna das curvas, a velocidade é
reduzida, favorecendo a deposição de sedimentos transportados pela água. Rios excessivamente sinuosos tendem a retardar
o escoamento, prolongando a permanência de água na bacia durante chuvas intensas e aumentando a probabilidade de
inundações. Por outro lado, a retificação artificial de meandros, embora acelere o fluxo, pode transferir riscos de
enchentes para áreas antes menos expostas.

\begin{figure}[htb]
	\caption{\label{sinus} Relação entre sinuosidade e velocidade em diferentes pontos de um rio}
	\begin{center}
	    \includegraphics[scale=0.40]{img/RiverSinus.png}
	\end{center}
	\legend{Fonte: Reprodução própria}
\end{figure} 

A carga de sedimentos é outra característica importante na morfologia dos rios. Ela reflete o volume de material
transportado, variando conforme a energia do fluxo e as propriedades do leito. Esta carga engloba tanto minerais
dissolvidos, derivados da erosão rochosa e responsáveis pela salinização de corpos receptores, quanto partículas sólidas
em suspensão, como fragmentos de rochas e matéria orgânica. O tamanho e o peso dessas partículas transportadas dependem
diretamente da velocidade do rio, dado que em fluxos mais lentos, as partículas tendem a se depositar no leito,
reduzindo sua mobilidade.

A descarga, por sua vez, representa o volume de água que passa por uma seção do rio em um determinado tempo e está
ligada à larguera e à profundidade do canal. A mesma pode ser dada em metros cúbicos de água por segundo ($m^3/s$) e é
definida pela da seção transversal do rio ($m^2$) e a velocidade da água ($m/s$) através de $D = A*v$, onde $D$ é a
descarga, $A$, a seção transversal e $v$ a velocidade. A descarga de um rio pode aumentar devido a fatores como a
precipitação no corpo d'água principal e em seus afluentes, pela contribuição de água subterrânea que se integra ao
sistema da bacia hidrográfica, ou pela quebra de barreiras naturais e/ou artificias, como barragens formadas por
sedimentos e matéria orgânica. Esses aumentos repentinos da descarga são os principais indicadores de enchentes, pois
podem superar a capacidade de armazenamento temporário do canal e das várzeas.

A relação entre descarga e enchentes é direta: sempre que a vazão excede a capacidade do canal, as águas transbordam
para as várzeas adjacentes e, em eventos extremos, para áreas além dessas planícies. Esses picos estão diretamente
associados à intensidade e duração das chuvas (\autoref{rain}), mas também ao estado prévio de umidade da bacia
hidrográfica. Em bacias já saturadas por chuvas anteriores, a água da precipitação escoa diretamente para os canais,
amplificando a descarga. Em contraste, solos secos absorvem grande parte da água, retardando o impacto no fluxo fluvial.

A capacidade de absorção do solo, determinada pela cobertura vegetal e pela porosidade – característica que determina a
taxa de infiltração da água e sua percolação\footnote{Movimento da água através das camadas do solo, favorecendo o
armazenamento em profundidade}, atua como um regulador natural. Solos porosos e ricos em matéria orgânica permitem maior
infiltração e percolação da água, enquanto compactação, urbanização ou camadas impermeáveis reduzem essa capacidade.
Quando o limite de saturação é atingido, o escoamento superficial direciona-se para canais menores, cuja sincronização
em múltiplos afluentes gera uma descarga acumulada crítica no rio principal, amplificando o risco de inundações. Logo,
um solo saturado e com pouca capacidade de inflitração resulta em enchentes mais frequentes em comparação com um solo
seco que absorve a água das chuvas sem exceder a capacidade dos canais.

A hidrodinâmica de enchentes fluviais resulta da interação complexa entre múltiplos fatores geomorfológicos e
hidrológicos. O transbordamento de corpos hídricos nunca decorre de um único elemento isolado, mas sim da interação
entre processos que, em condições críticas, potencializam eventos extremos. Podemos sintetizar esse mecanismo em três
estágios consecutivos: a ocorrência de precipitações intensas e prolongadas que excedem a capacidade de infiltração do
solo; o consequente escoamento superficial que eleva abruptamente a descarga hídrica; e a superação da capacidade de
vazão definida pela seção transversal do canal, levando à inundação das áreas adjacentes - particularmente as várzeas.
Esse processo evidencia como características físicas do canal, condições antecedentes do solo e padrões climáticos atuam
de forma integrada na gênese das enchentes.

% ######################################################################################################################
\subsection{Enchentes em ambientes urbanos e sistemas de prevenção}
Centros urbanos, conglomerados socioeconômicos críticos consolidados com frequente déficit de planejamento, operam
através de redes assimétricas de serviços essenciais (água, energia, transporte) cuja interdependência gera
vulnerabilidade sistêmica, com onde falhas pontuais levando à efeitos em cascata. Este quadro agrava-se diante das
projeções demográficas, pois, enquanto é estimado que a população humana aumente de 7.94 bilhões para 9.68 bilhões entre
o período de 2024 até 2050, (21\% em 26 anos), também é estimado que a população em centros urbanos cresça de 4.61
bilhões para 6.7 billhões durante 2024-2050, i.e., um crescimento de 30\% no mesmo período
\cite{World_Population_Prospects}. Tal expansão acelerada converterá áreas verdes (notadamente várzeas) em superfícies
impermeáveis, reduzindo a capacidade de infiltração do solo e ampliando o escoamento superficial pluvial.
Consequentemente, milhões de novos habitantes concentrar-se-ão em zonas com resiliência hidrológica comprometida,
expondo-os a riscos de enchentes superiores aos períodos pré-urbanização \cite{hammond}.

% \begin{figure}[htb] \caption{\label{katmandu} Enchente urbana em Kathmandu} \begin{center}
%   \includegraphics[scale=0.65]{img/kapan-flood-chandra.jpg} \end{center} \legend{Fonte: Chandra Bahadur Aale
%   \cite{chandra}} \end{figure}

As cidades enfrentam elevado risco de enchentes fluviais exatamente pelas alterações promovidas nessas zonas – comuns em
grandes centros urbanos devido à proximidade histórica de rios, dada a necessidade de acesso a água potável e descarte
de efluentes. Logo, as várzeas, outrora absorventes, tornam-se epicentros de risco. Quando transbordam, os rios não
apenas inundam essas planícies urbanizadas, mas também comprometem sistemas interdependentes, como sistemas de
distribuição de água, esgotos e infraestruturas de energia que sofrem danos em cascata, gerando impactos em regiões
distantes do epicentro do transbordamento.

A capacidade de eventos extremos causarem impactos significativos pode ser descrita através de três fatores: risco,
exposição e vulnerabilidade. Para que centros urbanos possam construir resiliência, é necessário primeiro entender estes
fatores. Desta forma, o planejamento urbano passa a ser menos reativo à situações de crise e mais proativo na sua
construção de adaptabilidade, levando à menores perdas de vidas em curto prazo e recuperação mais eficiente à longo
prazo, visando a preservação e restauração da funcionalidade de estruturas essênciais que compõem as malhas urbanas em
um ambiente em mudança. 

O risco, definido pela probabilidade de ocorrência do evento, está intrinsecamente vinculado aos fatores geográficos,
como padrões pluviométricos, e meteorológicos, como características topográficas, inerentes à localização do centro
urbano, assim como foi descrito na \autoref{morf}. Por depender destas variáveis naturais, trata-se do elemento sobre o
qual há menor controle direto por parte dos gestores urbanos, pois embora medidas governamentais possam ser tomadas
vizando a redução da crise climática, as mesmas são ineficazes quando aplicadas somente em escala regional.

% \begin{figure}[htb] \caption{\label{risco} Relação entre risco, exposição e vulnerabilidade em sistemas urbanos}
%   \begin{center} \includegraphics[scale=0.274]{img/RiverRisk.png} \end{center} \legend{Fonte: Bruno Merz et. al (2021)
%   \cite{merzcauses}} \end{figure}

A exposição, por outro lado, refere-se aos elementos materiais e econômicos suscetíveis a danos durante uma enchente,
sendo parcialmente controláveis por meio de políticas de planejamento urbano, como o zoneamento. Essa dimensão é
frequentemente quantificada pelo nível da água, que permite estimar quais ativos (infraestruturas, edificações, redes de
serviços) serão afetados em diferentes cenários de inundação. Estudos como o de Merz et al. \cite{merz} destacam,
entretanto, que essa métrica não pode explicar a totalidade dos danos, dado que fatores hidrodinâmicos adicionais, como
velocidade da corrente e carga de sedimentos, desempenham papéis críticos.

Apesar dessas nuances, o nível da água mantém utilidade como indicador generalizado de impacto, dado que os outros
fatores são altamente dependentes do tipo de estrutura afetado, e, logo, o mesmo oferece uma base objetiva para
estratégias de zoneamento urbano que visem minimizar perdas. Além disso, o monitoramento do nível da água representa uma
forma unificada de entender dinâmicas complexas de inundação. Embora o conhecimento do mesmo só seja efetivo caso
aplicado em mapas de inundação, a redução da complexidade para um único valor facilita a comunicação do risco para a
população e agiliza a tomada de decisão baseado nas capacidade das estruturas de defesa presentes. Neste contexto,
sistemas de monitoramento contínuo do nível d'água assumem papel central, pois fornecem dados em tempo real que
alimentam modelos preditivos e sistemas de alerta precoce.

A vulnerabilidade, por sua vez, é o fator de maior controle para tomadores de decisões de centros urbanos. Estas são as
medidas de adaptabilidade tanto de estruturas, quanto de habitações perante eventos extremos. As medidas relacionadas à
redução da vulnerabilidade incluem desde infraestruturas físicas, como sistemas de defesa contra as enchentes, sistemas
de prevenção, até sistemas de gerenciamento de crises e alerta de riscos. 

Em escala regional, barragens destacam-se como sistemas estratégicos de armazenamento hídrico, capazes de regular fluxos
extremos provenientes de chuvas intensas. Além da geração de energia, as represas podem alocar partes dos seus
reservatórios para contenção de cheias, liberando água de forma gradual no curso principal da bacia hidrográfica de
acordo com a capacidade da mesma. Essa modulação reduz a velocidade de escoamento, minimizando danos por erosão e
transbordamentos.

% \begin{figure}[htb] \caption{\label{itaipu} Usina hidroelétrica de Itaipu com seus vertedouros abertos} \begin{center}
%   \includegraphics[scale=0.12]{img/Itaipu_geral.jpg} \end{center} \legend{Fonte: Jonas de Carvalho \cite{itaipu}}
%   \end{figure}

Uma solução mais comum em centros urbanos se dá nas formas do redirecionamento dos cursos naturais dos rios, em especial
com a retificação de canais naturais substituindo a sinuosidade fluvial por trajetos retilíneos artificialmente
construídos, dado que canais sinuosos diminuem a velocidade as águas em seus meandros e retardam o escoamento (vide
\autoref{sinus}). Os canais retificados não apenas direcionam a vazão crítica para corpos receptores com menor risco de
transbordamento, mas também integram-se à infraestrutura urbana como componentes de sistemas de abastecimento hídrico e
monitoramento de condições fluviais. Além disso, eles podem contar com sistemas de desvio que visam ligar partes
distintas da rede hidrográfica, desviando parcelas da vazão para rotas alternativas que contornam núcleos urbanos
densamente povoados.

Como alternativa complementar à retificação de canais, o aterramento fluvial amplia a capacidade de contenção hídrica
mediante estruturas de expansão das margens – como diques ou aterros. Essas intervenções aumentam a seção transversal do
canal, aumentando sua capacidade de vazão antes do transbordamento. Contudo, a eficácia de ambas estas estruturas
depende criticamente da integridade estrutural: rupturas pontuais, como brechas em diques ou falhas de vedação,
convertem essas barreiras em focos de amplificação hidráulica. Nessas situações, a súbita liberação de energia cinética
concentrada em áreas adjacentes, muitas vezes desprovidas de drenagem redundante e com falta de aviso prévio, exacerba
inundações locais.

Além das intervenções superficiais, sistemas subterrâneos de drenagem pluvial complementam a gestão hidráulica urbana.
Projetadas para transportar água separadamente do esgoto, as galerias pluviais operam como redes de canais subterrâneos
que direciona o escoamento superficial para canais, sejam eles artificiais ou naturais, ou para estruturas de contenção.

% -> Essas galerias pluviais podem até mesmo apresentar riscos maiores quando levado em consideração a sua habitação por
% moradores de rua, como nos túneis de Las Vegas. Em situações de chuvas intensas, nas quais os túneis serviriam o seu
% propósito de escoamento alternativo, essa parcela da população se encontra em um enorme risco de afogamento, visto que
% o fluxo da água nessas estruturas é súbito e intenso.

Integradas a essas redes, estações de bombeamento hidráulico atuam na transferência controlada de água entre bacias,
equalizando cargas em áreas de gradiente topográfico desfavorável. Em escala subterrânea, reservatórios de amortecimento
operam como análogos submersos de lagos artificiais, armazenando volumes críticos em períodos de pico para posterior
liberação gradual via sistemas de drenagem. Esses reservatórios permitem tanto a descarga regulada para canais
artificiais quanto a inflitração direta do lençol freático, contornando a impermeabilidade do solo urbano através destes
processos de percolação assistida. Essa dualidade funcional transforma a infraestrutura em mecanismo híbrido de controle
e restauração hidrológica.

Os reservatórios de retenção seca representam soluções multifuncional que operam como espaços públicos ordinários e
bacias de acumulação durante eventos extremos. Localizados estrategicamente em áreas de baixa cota topográfica (como
parques e praças), aproveitam o gradiente natural para interceptar fluxos superficiais, reduzindo a carga sobre redes de
drenagem em períodos chuvosos enquanto mantêm funções recreativas em condições normais. Sua integração urbanística
minimiza impactos visuais e sociais, com zonas de infraestrutura de baixa criticalidade que garantem rápida recuperação
operacional pós-inundações - característica particularmente valiosa em áreas densamente ocupadas com espaço limitado. O
Japão, devido à sua alta densidade urbana, exemplifica esta abordagem no complexo Parque Shin-Yokohama/Estádio Nissan
\cite{shinyokohamaplan}, projetado para conter enchentes do rio Tsurumi adjacente (\autoref{shin_yoko}). Durante cheias,
a água é direcionada ao parque por diques (\autoref{shin_yoko}.1), armazenada temporariamente (\autoref{shin_yoko}.2) e
liberada controladamente quando o canal recupera capacidade (\autoref{shin_yoko}.3), complementado por alterações no
canal fluvial e elevação estrutural do estádio para garantir funcionalidade durante eventos extremos.

\begin{figure}[htb]
	\caption{\label{shin_yoko} Parque de Shin-Yokohama}
	\begin{center}
	    \includegraphics[scale=2.3]{img/Lago-parque-japao.png}
	\end{center}
	\legend{Fonte: Governo da Prefeitura de Kanagawa \cite{shinyokohamapark} e Ministério de Terras, Infraestrutura,
	Transporte e Turismo, Escritório de Desenvolvimento Regional de Kanto \cite{shinyokohamaplan}}
\end{figure}


% Shin-Yokohama Park (新横浜公園): https://www.nissan-stadium.jp/csr/safety03.php Estádio Nissan (日産スタジアム):
% https://nissan-stadium.j-server.com/LUCNISSANS/ns/tl.cgi/https://www.nissan-stadium.jp/csr/safety03.php
% https://www.ktr.mlit.go.jp/keihin/keihin00119.html
% https://www-ktr-mlit-go-jp.translate.goog/keihin/keihin00119.html?_x_tr_sl=auto&_x_tr_tl=en&_x_tr_hl=en-US

Enquanto as Infraestruturas Cinza, grandes sistemas de engenharia convencional em concreto para controle direto de
fluxos, representam soluções predominantes, as Infraestruturas Verdes emergem como abordagem complementar baseada em
processos naturais. Estas utilizam componentes vegetados e permeáveis (parques, jardins, telhados vivos e rearborização
de vias) para substituir superfícies impermeáveis, aumentando a capacidade de absorção urbana. Sua implementação em
zonas industriais, comerciais e residenciais visa não apenas a gestão sustentável de águas pluviais, mas também a
mitigação de desequilíbrios urbanos, proporcionando melhorias da qualidade do ar, redução de ilhas de calor e incremento
da biodiversidade. Esta multifuncionalidade torna-as particularmente valiosas em áreas de alta densidade, contrastando
com as limitações das estruturas cinza: elevados custos de manutenção, potencial deslocamento de riscos para áreas
adjacentes e menor resiliência ecológica. Quando integradas como sistemas complementares, as soluções verdes oferecem
suporte sustentável à redução de riscos de enchentes urbanas \cite{gill}.

% A mitigação de riscos hidrológicos urbanos exige a complementaridade entre medidas estruturais e não estruturais,
% sendo estas últimas fundamentais para reduzir a vulnerabilidade sistêmica. Medidas de zoneamento emergem como
% estratégia central, delimitando restrições à ocupação em planícies aluviais mediante modelagem hidrodinâmica de
% cenários de enchente. A relocação preventiva de infraestruturas críticas nessas zonas – embora politicamente complexa
% – reduz exponencialmente a exposição de ativos estratégicos, enquanto a reconversão de áreas ocupadas em zonas verdes,
% como parques ou áreas públicas. Além disso, para infraestruturas que não podem ser removidas destas áreas, é essencial
% a adaptação das mesmas buscando mais resiliência.

A eficácia dessas ações depende criticamente de sistemas preditivos integrados, combinando redes de monitoramento
hidrometeorológico, como estações fluviométricas e pluviométricas para monitoramento de níveis de rios e eventos
extremos, com plataformas de análise em tempo real. Tais sistemas permitem a geração de alertas precoces baseados em
limiares de precipitação acumulada e níveis hidrométricos críticos. Contudo, o ciclo de gestão de risco só se completa
com protocolos de comunicação efetiva, que convertam estes dados técnicos em informações acionáveis para a população.
Sistemas de alerta georreferenciados, como a implementação de sirenes moduladas, avisos via mensagens de SMS e APIs para
integração em redes sociais, devem ser continuamente adaptáveis para garantir a comunicação efetiva dos riscos e
aumentar a segurança comunitária.

Um exemplo de sistema de monitoramento integrado nacional vem através do Centro Nacional de Monitoramento e Alertas de
Desastres Naturais (CEMADEN) que opera uma rede de monitoramento hidrometeorológico no Brasil. Seu modelo combina coleta
de dados in situ através de sensores pluviométricos e fluviométricos instalados em áreas de risco, com análises e
imagens de satélite \cite{cemadem}. Essa infraestrutura permite a emissão de alertas antecipados para eventos extremos
(como enchentes e deslizamentos), direcionados a defesas civis municipais e estaduais. Contudo, desafios operacionais
persistem, incluindo falhas na atualização de limiares críticos e lacunas na densidade de cobertura instrumental, como
evidenciado durante as enchentes de 2024 no Rio Grande do Sul, onde limitações na transmissão de alertas impactaram a
eficácia das respostas emergenciais.